Choose an affine open $U\subseteq X$ meeting $Y$, and let $\mathfrak p=I(\overline Y\cap U)\subseteq A(U)$. A rational function regular on an open set meeting $Y$ belongs to $A(U)_{\mathfrak p}$: choose a point of that intersection and a local quotient whose denominator is nonzero there. Conversely, a denominator outside $\mathfrak p$ is nonzero on a nonempty open subset of $Y$. Hence $\mathcal O_{Y,X}=A(U)_{\mathfrak p}$.

This is local with maximal ideal $\mathfrak p A(U)_{\mathfrak p}$, residue field $\operatorname{Frac}(A(U)/\mathfrak p)=K(Y)$, and dimension $\operatorname{ht}\mathfrak p=\dim X-\dim Y$ by (1.8A). If $\dim Y>0$, choose a transcendence basis $t_1,\ldots,t_r$ of $K(Y)/k$. The extension over $k(t_1,\ldots,t_r)$ is finite. For arbitrarily large integers $m$ prime to the characteristic, $T^m-t_1$ is irreducible over that rational function field by Eisenstein at $t_1$. An algebraically closed $K(Y)$ would contain a root of each, contradicting the finite extension degree. Thus $K(Y)$ is not algebraically closed.
